\section{Optimized Inference I/O for Full Attention Residuals}
\label{app:inference}

A na\"ive implementation of Full AttnRes scans all preceding layer outputs at every layer, so memory traffic scales linearly with depth. As noted in \S\ref{sec:infra-inference}, however, the pseudo-query $\bm{w}_l$ is a learned parameter independent of both the input and the hidden state. We can therefore batch inter-block accesses across layers in a two-phase schedule, bringing total I/O well below the na\"ive bound.

Note that the block partition introduced below is purely an inference scheduling device. Unlike Block AttnRes, it leaves the model architecture unchanged and does not replace per-layer sources with block summaries; it simply makes the amortization argument concrete.

\paragraph{Setup}

Let the model have $L$ layers and hidden dimension $d$, partitioned into $N$ contiguous blocks of size $S = L/N$. Inference proceeds one block at a time: Phase~1 jointly computes inter-block attention for all $S$ layers in the block against all preceding blocks, and Phase~2 walks through intra-block dependencies sequentially.

\subsection*{Phase 1: Batched Inter-block Attention}

Consider block $n$ with its $S$ layers. The queries $\{\bm{w}_l\}_{l \in \mathcal{B}_n}$ are all known before execution begins, so the $(n{-}1)S$ preceding key--value pairs need only be read once from HBM and reused across all $S$ queries. The read cost for block $n$ is therefore
\begin{equation}
    \mathrm{Read}_{\text{inter}}^{(n)} = 2(n-1)Sd,
\end{equation}
where the factor of $2$ accounts for both keys and values. Summing over all $N$ blocks and using $SN=L$:
\begin{equation}
    \mathrm{Read}_{\text{inter}} = \sum_{n=1}^{N} 2(n-1)Sd
    = 2Sd \cdot \frac{N(N-1)}{2}
    = dL(N-1).
\end{equation}

Phase~1 also writes one $d$-dimensional output per layer, giving $\mathrm{Write}_{\text{inter}}^{(n)} = Sd$ per block and
\begin{equation}
    \mathrm{Write}_{\text{inter}} = Ld
\end{equation}
in total.

\subsection*{Phase 2: Sequential Intra-block Attention}

Phase~1 covers all sources before the current block. Within the block, however, each layer depends on those before it, so these must be handled in order. Layer $t$ ($1 \le t \le S$) reads $t{-}1$ intra-block key--value pairs at a cost of $2(t{-}1)d$. Summing over one block:
\begin{equation}
    \mathrm{Read}_{\text{intra}}^{(n)} = \sum_{t=1}^{S} 2(t-1)d = S(S-1)d.
\end{equation}
Phase~2 also writes one output per layer, so $\mathrm{Write}_{\text{intra}}^{(n)} = Sd$.

\subsection*{Total Amortized I/O per Layer}

Summing both phases over all $N$ blocks:
\begin{equation}
    \mathrm{Read}_{\text{total}} = dL(N-1) + N \cdot S(S-1)d, \qquad
    \mathrm{Write}_{\text{total}} = 2Ld.
\end{equation}
Dividing by $L$ and using $SN=L$:
\begin{equation}
    \text{Read per layer} = (N-1)d + (S-1)d = (S + N - 2)d, \qquad
    \text{Write per layer} = 2d,
\end{equation}
\begin{equation}
    \boxed{\;\text{Total I/O per layer} = (S + N)\,d.\;}
\end{equation}

Batching inter-block reads thus brings per-layer I/O from $\mathcal{O}(L)$ down to $\mathcal{O}(S{+}N)$. The schedule follows the same two-phase split as Block AttnRes: inter-block attention accounts for the bulk of the traffic, while sequential computation stays local within each block.
